\documentclass[11pt]{article}
\usepackage[margin=1in]{geometry}
\usepackage{amsmath,amssymb,amsthm}
\usepackage{enumitem}

\newtheorem{theorem}{Theorem}
\newtheorem{lemma}[theorem]{Lemma}
\newtheorem{proposition}[theorem]{Proposition}
\newtheorem{corollary}[theorem]{Corollary}
\newtheorem{definition}[theorem]{Definition}
\newtheorem{remark}[theorem]{Remark}

\newcommand{\Pair}{\mathrm{Pair}}
\newcommand{\TreeEq}{\mathrm{TreeEq}}
\newcommand{\IfLf}{\mathrm{IfLf}}
\newcommand{\Fst}{\mathrm{Fst}}
\newcommand{\Snd}{\mathrm{Snd}}
\newcommand{\KT}{\mathrm{KT}}
\newcommand{\Comp}{\mathrm{Comp2}}
\newcommand{\Rec}{\mathrm{Rec}}
\newcommand{\SchemaF}{\textbf{Schema~F}}
\newcommand{\Deriv}{\mathrm{Deriv}}
\newcommand{\poo}{\Pair(O,O)}
\newcommand{\Fun}{\mathrm{Fun1}}
\newcommand{\geval}{\mathrm{geval}}
\newcommand{\thFun}{\mathrm{thFun}}
\newcommand{\reify}{\mathrm{reify}}
\newcommand{\code}{\mathrm{code}}
\newcommand{\iterf}{\mathrm{iterate}}

\title{The Boundary of Nelson's Feasibility Program:\\
  Per-Program Proofs Are Feasible,\\
  Self-Reference Is Not}

\author{}
\date{}

\begin{document}
\maketitle

\begin{abstract}
We formalize in Agda a tree-based equational system with Schema~F
(tree induction) and analyze the feasibility of consistency proofs
in the sense of Nelson's program.  The formalization contains:
(1)~a complete diagonal G\"odel~II proof (168 lines, 0 postulates);
(2)~a micro-Nelson instance verifying 56 programs against a fixed target;
(3)~a uniform proof-construction mechanism (the R/Q technique) that produces
constant-size, BLevel-0 invariant proofs for all oblivious step functions;
(4)~the Kritchman--Raz contradiction theorem (\texttt{krClean}), fully proved.
The main conclusion: the \emph{per-program} content of a consistency proof
is feasible (O(1), ground-verifiable), but combining these proofs
into a system-level consistency statement requires self-reference,
which is irreducibly infeasible.  The boundary is precisely located:
it is the transition from ``for each program~$p$'' to ``for all
programs of size~$\le l$.''
\end{abstract}

%% ---------------------------------------------------------------
\section{The equational system}

The system operates on binary trees:
\[
  t \;::=\; \mathrm{lf} \mid \mathrm{nd}(t_1, t_2)
\]
with $O = \mathrm{lf}$ and $\Pair(a,b) = \mathrm{nd}(a,b)$.
Terms are built from variables, constants~$O$, and applications of
one-argument functors~$\Fun$ and two-argument functors~$\mathrm{Fun2}$:
\[
  \Fun: \quad I,\; \Fst,\; \Snd,\; \Comp\; h\; f\; g,\;
    \Rec\; z\; s,\; \KT\; c
\]
\[
  \mathrm{Fun2}: \quad \Pair,\; \IfLf,\; \TreeEq,\;
    \mathrm{Lift}\; f,\; \mathrm{Post}\; f\; h,\; \mathrm{Fan}\; h_1\; h_2\; h
\]
Equations are pairs of terms.  Derivations $\Deriv\; \mathit{hyp}\; (l = r)$
are built from computation axioms ($\mathrm{axI}$, $\mathrm{axFst}$,
$\mathrm{axSnd}$, $\mathrm{axComp2}$, $\mathrm{axRecLf}$,
$\mathrm{axRecNd}$, $\mathrm{axIfLfL}$, $\mathrm{axIfLfN}$,
$\mathrm{axTreeEqLL}$, \ldots), structural rules
($\mathrm{ruleTrans}$, $\mathrm{ruleSym}$, $\mathrm{congL}$,
$\mathrm{congR}$, $\mathrm{cong1}$), substitution ($\mathrm{ruleInst}$),
and \SchemaF{} ($\mathrm{ruleF}$).

\SchemaF{} is tree induction: if $f$ and $g$ agree on $O$ and
satisfy the same step equation with respect to a step functor~$s$,
then $f(x) = g(x)$ for all~$x$.
Consistency means no derivation of $O = \poo$.

%% ---------------------------------------------------------------
\section{Proof encoding and ground evaluation}

The function $\thFun : \mathrm{Tree} \to \mathrm{Tree}$ interprets
a tree as a proof encoding: given a tree~$t$, it produces
$\mathrm{nd}(\mathit{lhs}, \mathit{rhs})$ representing the equation
that~$t$ encodes.  There are 26 proof-encoding tags
(axioms~0--16, structural rules~17--25).

The ground evaluator $\geval$ dispatches on functor-code tags
(offset~26--32) and computes the semantic value of coded ground terms.
The proof checker $\mathrm{checkProof2}(t)$ runs $\thFun$ to extract
the equation, then checks whether $\geval$ agrees on both sides.

Key property: for \emph{ground} coded terms (no variables),
$\geval$ correctly computes the semantic value.  When variables
appear, $\geval$ returns a stuck term.  The distinction between
ground and stuck is exactly the distinction between oblivious and
non-oblivious proofs.

%% ---------------------------------------------------------------
\section{Programs and iteration}

A \emph{program} is a $\Fun$ of the form $\iterf\; f\; O$,
where $f$ is a step function and $O$ is the initial value.
The combinator $\iterf\; f\; O = \Rec\; O\; (\mathrm{iterStep}\; f)$
satisfies:
\begin{align*}
  \iterf\; f\; O\; (O) &= O \\
  \iterf\; f\; O\; (\Pair(a,b)) &= f(\iterf\; f\; O\; (b))
\end{align*}
Applied to $\mathrm{natCode}(n) = \mathrm{nd}(\mathrm{lf},
\mathrm{natCode}(n{-}1))$, this gives $f^n(O)$.

The \emph{code size} of a program is
$\mathrm{treeSize}(\mathrm{codeF1}(p))$, counting nodes
in the tree representation of the functor code.

%% ---------------------------------------------------------------
\section{The R/Q technique}

Two lemmas, both proved by \SchemaF{} at BLevel~0:

\begin{lemma}[R --- ifLfCollapse]
  For all terms $x$:
  \[
    \IfLf(x, \Pair(\poo, \poo)) = \poo
  \]
\end{lemma}
\begin{proof}
  Define $R(x) = \IfLf(x, \Pair(\poo, \poo))$.
  Base: $R(O) = \poo$ by $\mathrm{axIfLfL}$.
  Step: $R(\Pair(v_0, v_1)) = \poo$ by $\mathrm{axIfLfN}$.
  Both match $\KT(\poo)$.  \SchemaF{} gives $R(x) = \poo$ for all~$x$.
\end{proof}

\begin{lemma}[Q --- for $a \ne b$ with one component $O$]
  Let $a = \poo$, $b = O$.  For all terms~$x$:
  \[
    \IfLf(\TreeEq(x, a), \Pair(\TreeEq(x, b), \poo)) = \poo
  \]
\end{lemma}
\begin{proof}
  Define $Q(x) = \IfLf(\TreeEq(x, \poo), \Pair(\TreeEq(x, O), \poo))$.
  Base: $Q(O)$.  $\TreeEq(O, \poo) = \poo$ by $\mathrm{axTreeEqLN}$.
  $\IfLf(\poo, \ldots) = \poo$ by $\mathrm{axIfLfN}$.
  Step: $Q(\Pair(v_0, v_1))$.  $\TreeEq(\Pair(v_0,v_1), O) = \poo$
  by $\mathrm{axTreeEqNL}$.  So the $\Pair$ argument becomes
  $\Pair(\poo, \poo)$.  Apply~R.  \SchemaF{} closes.
\end{proof}

The key: Q captures the case analysis that would normally require
knowing the value of~$x$.  When $x = a$, the first comparison gives~$O$
and falls through to the second, which detects $x \ne b$ (since $a \ne b$).
When $x \ne a$, the first comparison gives $\poo$ and short-circuits.
\SchemaF{} handles this without case-splitting because both branches
produce~$\poo$.

%% ---------------------------------------------------------------
\section{Three proof families}

For step functions $f = \Comp\;\Pair\; g\; h$ where
$g, h \in \{I, \Fst, \Snd, \KT\; O, \KT\; \poo\}$, and target
$\xi_1 = \Pair(\poo, O)$, three parametric proof constructors
handle all oblivious cases.

\subsection*{Family A: Constant $g$, Fst mismatch}
If $g = \KT\; c$ with $\TreeEq(c, \Fst(\xi_1)) = \poo$
(i.e., $c \ne \poo$), then in the step case:
\[
  \TreeEq(\Pair(c, h(\mathrm{rec}_1)), \xi_1)
  = \IfLf(\underbrace{\TreeEq(c, \poo)}_{= \poo}, \ldots)
  = \poo
\]
The Fst comparison short-circuits.  The h~side is never examined.
\textbf{5 programs} (all with $g = \KT\; O$).

\subsection*{Family B: Non-constant $g$, constant $h = \KT\;\poo$}
The Snd comparison gives $\TreeEq(\poo, O) = \poo$ by
$\mathrm{axTreeEqNL}$.  The IfLf argument becomes
$\Pair(\poo, \poo)$.  R~collapses it.
\textbf{3 programs} ($g \in \{I, \Fst, \Snd\}$).

\subsection*{Family C: $g = h$ (same function), Q lemma}
$f(x) = \Pair(g(x), g(x))$.  The comparison with
$\xi_1 = \Pair(\poo, O)$ gives:
\[
  \IfLf(\TreeEq(g(\mathrm{rec}_1), \poo),\;
        \Pair(\TreeEq(g(\mathrm{rec}_1), O), \poo))
  = Q(g(\mathrm{rec}_1)) = \poo
\]
\textbf{3 programs} ($g = h \in \{I, \Fst, \Snd\}$).

\medskip
\noindent\textbf{Total:} 11 programs with full equational ($\Deriv$) proofs.
All 23 Comp2 Pair combinations pass $\mathrm{checkProof2}$ (geval verification).

%% ---------------------------------------------------------------
\section{The micro-Nelson instance}

\texttt{MicroNelson.agda} enumerates 28 step functions
(all $\Comp\;\Pair\; g\; h$ with $g, h \in \{I, \Fst, \Snd,
\KT\;O, \KT\;\poo\}$, plus 5~base $\Fun$) in two modes
(iterate + direct), giving 56 programs.  The target is
$\xi_4 = \Pair(\Pair(\Pair(O,O),O), \Pair(O,O))$, chosen deeper
than any $f(O)$ for depth-1 step functions.

All 56 satisfy $\mathrm{checkProof2}(\mathrm{mkTest}(p, \xi_4)) = \mathrm{true}$,
verified by Agda \texttt{refl}.  The adversarial program
$\KT(\reify\;\xi_4)$ (which always outputs~$\xi_4$) correctly
returns \texttt{false}.

This is a complete pigeonhole instance at depth~$\le 1$:
every program of bounded constructor-depth is accounted for,
and $\xi_4$ is not in any program's output.

%% ---------------------------------------------------------------
\section{The Kritchman--Raz framework}

The Chaitin machine $C_l$ searches for proofs that individual
programs don't output some string~$\xi$.  If the system proves
its own consistency, the search terminates and $C_l$ outputs~$\xi$.
But $|C_l|$ is bounded (the step function is fixed; $l$ enters
only as data via binary encoding, contributing $O(\log l)$ to
the code size).  For large~$l$, $|C_l| < l$, so $C_l$ itself
witnesses $K(\xi) \le |C_l| < l$, contradicting $K(\xi) > l$.

\begin{theorem}[\texttt{krClean}, fully proved in Agda]
  Given:
  \begin{enumerate}[nosep]
    \item A Chaitin machine family $\mathit{cm} : \mathbb{N} \to \Fun$,
    \item A threshold $T$ with $|\mathit{cm}(l)| \le l$ for $l \ge T$,
    \item For each $l \ge T$: a tree $\xi$, a clock $c$,
      a proof that $\mathit{cm}(l)$ outputs $\xi$ at clock~$c$,
      and a proof that no $\Fun$ of size $\le l$ outputs~$\xi$,
  \end{enumerate}
  then $\bot$.
\end{theorem}

The proof is direct: $\mathit{cm}(T)$ outputs $\xi$ at clock~$c$
and has $|\mathit{cm}(T)| \le T$, contradicting the incompressibility
of~$\xi$.

%% ---------------------------------------------------------------
\section{The boundary}

The formalization reveals a sharp decomposition of the consistency proof:

\begin{center}
\begin{tabular}{lll}
\textbf{Component} & \textbf{Status} & \textbf{Obstacle} \\
\hline
Per-program invariant & Feasible & None (O(1), BLevel 0) \\
Ground verification & Feasible & None (checkProof2, refl) \\
Pigeonhole (finite $l$) & Feasible & None (enumeration) \\
\hline
$\forall p.\; |p| \le l \Rightarrow \ldots$ & Infeasible & Self-reference \\
Chaitin machine as $\Fun$ & Infeasible & Self-reference \\
$|C_l| < l$ for large $l$ & Infeasible & Self-reference \\
\end{tabular}
\end{center}

The transition from feasible to infeasible is precisely the
transition from ``for each specific program~$p$'' to
``for all programs of size~$\le l$.''

\begin{itemize}[nosep]
\item ``For each~$p$'': the R/Q technique produces a constant-size
  proof.  The proof is oblivious --- it doesn't examine what $p$
  computes, only the structural mismatch between $p$'s output
  format and~$\xi$.

\item ``For all~$p$'': requires the system to enumerate its own
  programs, count them (pigeonhole), and recognize that the Chaitin
  search procedure is itself a program (self-reference).  This is
  the same obstacle as in the diagonal proof, repackaged.
\end{itemize}

The diagonal proof (GodelII.agda) concentrates self-reference in
the substitution machinery ($\mathrm{substTFor}$, $\mathrm{codedSubst}$,
variable capture).  The KR approach distributes it across counting
and search.  Both are blocked by the same fundamental constraint:
the system cannot fully reason about its own proof structure.

%% ---------------------------------------------------------------
\section{Implications for Nelson's program}

\begin{enumerate}
\item \textbf{The per-program feasibility claim holds.}
  Each invariant proof is O(1), BLevel~0, and ground-verifiable.
  The R/Q technique provides this uniformly for all oblivious
  programs.  The $\geval$ evaluator extends coverage to
  non-oblivious programs at the metalevel.  This is the positive
  content of Nelson's program.

\item \textbf{The organizational infeasibility persists.}
  Combining per-program proofs into a system-level consistency
  statement requires self-reference --- whether via diagonalization
  or via counting.  The KR approach, which avoids the diagonal
  lemma entirely, hits the same wall.  The obstacle is not an
  artifact of the proof method.

\item \textbf{The boundary is intrinsic.}
  The oblivious/non-oblivious distinction from AutoProof is a local
  version of the each/all boundary.  Oblivious proofs work by
  structure alone (the output format mismatches~$\xi$).
  Non-oblivious proofs require knowing the actual computed value ---
  a miniature version of the self-reference problem.

\item \textbf{Bounded consistency is feasible.}
  For a fixed depth bound~$d$, the micro-Nelson instance gives
  a complete, computationally verified consistency proof:
  every program of depth~$\le d$ is checked against~$\xi$,
  and all pass.  This is a concrete, finite, feasible consistency
  result.  It does not extend to unbounded~$d$ without
  self-reference.
\end{enumerate}

%% ---------------------------------------------------------------
\section{Formalization summary}

All results are formalized in Agda 2.9.0 with
\texttt{-{}-safe -{}-without-K -{}-exact-split}.  Zero postulates throughout.

\medskip
\begin{center}
\begin{tabular}{lrl}
\textbf{File} & \textbf{Lines} & \textbf{Content} \\
\hline
GodelII.agda & 168 & Diagonal G\"odel~II \\
MicroNelson.agda & 258 & 56 programs, checkProof2 = true \\
AutoProof.agda & 754 & R/Q library, 3 proof families, 11 Deriv proofs \\
Chaitin.agda & 395 & Binary encoding, tree enumeration, krClean \\
KRShortProof.agda & 310 & Case~1 template (Fst mismatch) \\
DoublerProof.agda & 430 & Case~3 template (Q lemma) \\
GeneralQ.agda & 443 & Cases~2,4; tripler, mixed \\
SwapProof.agda & 492 & Case~4 template (h2 helper) \\
\hline
Guard/ total & $\sim$40 files & Full system including infrastructure \\
\end{tabular}
\end{center}

\end{document}
